\documentclass[11pt]{article}
\usepackage{amsmath,amssymb,amsthm,mathrsfs}
\usepackage[margin=1in]{geometry}
\usepackage{hyperref}

\newtheorem{theorem}{Theorem}[section]
\newtheorem{lemma}[theorem]{Lemma}
\newtheorem{proposition}[theorem]{Proposition}
\newtheorem{corollary}[theorem]{Corollary}
\theoremstyle{definition}
\newtheorem{definition}[theorem]{Definition}
\newtheorem{remark}[theorem]{Remark}
\newtheorem{example}[theorem]{Example}
\newtheorem{question}[theorem]{Question}
\newtheorem{conjecture}[theorem]{Conjecture}
\newtheorem{heuristic}[theorem]{Heuristic}

\newcommand{\N}{\mathbf{N}}
\newcommand{\Nz}{\mathbf{N}_0}
\newcommand{\Z}{\mathbf{Z}}
\newcommand{\Q}{\mathbf{Q}}
\newcommand{\R}{\mathbf{R}}
\newcommand{\C}{\mathbf{C}}
\newcommand{\F}{\mathbf{F}}
\newcommand{\SLNz}{SL_2(\Nz)}
\newcommand{\SLN}{SL_2(\N)}
\newcommand{\Df}{\mathcal{D}_f}
\newcommand{\Dphi}{\mathcal{D}_{\phi_0}}
\newcommand{\Sb}{\bar{S}}
\newcommand{\Tb}{\bar{T}_f}
\newcommand{\cb}{\bar{c}_f}
\newcommand{\Ff}{\hat{F}_f}
\newcommand{\Phzero}{\Phi_0}
\newcommand{\PP}{\mathbf{P}}
\newcommand{\Iset}{\mathcal{I}_{\phi_0}}
\newcommand{\Spine}{\mathfrak{S}}
\newcommand{\depth}{\mathrm{depth}}
\newcommand{\norm}[1]{N(#1)}
\newcommand{\Zi}{\Z[i]}

\title{Combinatorics of the interior-spine sieve: \\ generating functions, the Sundaram dictionary, and depth-bounded densities}
\author{(working note in the Shakov framework, P2)}
\date{April 2026}

\begin{document}
\maketitle

\begin{abstract}
We develop the interior-spine sieve of \cite{B8} as a precise combinatorial object. Working primarily with $f = \phi_0(n) = n^2+1$, we prove an exact functional equation for the trivariate generating function
\[
G(x,y,z) := \sum_{A \in \SLNz} x^{\chi(A)} y^{m(A)} z^{m(c(A))}
\]
which packages the divisor pairs of $\Phzero$ together with their complements under the involution. Specialising to $z=1, y=1$ recovers the divisor-twisted generating function $\sum_{n\ge 0} \tau(n^2+1) x^n$. We then establish a precise dictionary between Sundaram's classical sieve and the interior-spine sieve: each Sundaram exclusion is a depth-$1$ spine exclusion in a degenerate (non-enumerable) ambient monoid, and conversely the interior-spine sieve reduces to Sundaram-type APs but enriched by Gaussian-integer witnesses. We quantify the surplus information by exhibiting concrete families of composite indices that the spine sieve detects which Sundaram does not. We compute exact densities of spine exclusion at depths $\le 2$ and $\le 3$, prove an upper bound at depth $\le k$ in terms of singular series, and give a multiplicity-counting identity that reproduces the divisor function. We extend the generalised lemma of \cite{B8} by giving a sharp \emph{second-tier} criterion. Throughout we are careful to flag every result that is conjectural: the principal new theorems are (\ref{thm:fnceq}), (\ref{thm:multiplicity}), (\ref{thm:sundaram-dict}), (\ref{thm:depth2-density}), (\ref{thm:depth-k-upper}), and (\ref{thm:second-tier}); the principal conjectures are (\ref{conj:depth-asymp}), (\ref{conj:joint-density}). No claim of progress on Landau's fourth problem is made.
\end{abstract}

\tableofcontents

\section{Setting and notation}

We use the conventions of Shakov \cite{shakov} and the working note \cite{B8}. The free monoid $\SLNz$ is generated by
\[
S = \begin{pmatrix}1 & 0 \\ 1 & 1\end{pmatrix}, \qquad T = \begin{pmatrix}1 & 1 \\ 0 & 1\end{pmatrix},
\]
with complement involution $c(A) = JAJ$ for $J = \bigl(\begin{smallmatrix}0 & 1 \\ 1 & 0\end{smallmatrix}\bigr)$ swapping $S \leftrightarrow T$. The interior $\SLN \subset \SLNz$ is the set of matrices with all four entries strictly positive; equivalently the set of $S/T$-words containing both letters at least once.

For $f \in \{\phi_0, \phi_1, \psi_2, \phi_3\}$ enumerable in the sense of \cite[Defn.~2.1]{shakov}, $\Ff: \SLNz \to \Df$ is the unique invertible equivariant map sending $I \mapsto (1,0)$. We focus throughout on $f = \phi_0(n) = n^2+1$ and the map
\[
\Phzero: \SLNz \to \Dphi, \qquad \Phzero \begin{pmatrix} a & b \\ c & d \end{pmatrix} = \bigl(a^2+b^2,\ ac+bd\bigr).
\]
Write $m(A) := a^2+b^2$ and $\chi(A) := ac+bd$. Note $m(c(A)) = c^2+d^2$ is the bottom-row norm, equivalently the top-row norm of the complement. Recall the underlying Diophantus identity
\[
\chi(A)^2 + (ad-bc)^2 = m(A) \cdot m(c(A)),
\]
which on $\SLNz$ (where $ad-bc = 1$) reduces to
\begin{equation}\label{eq:diophantus}
\chi(A)^2 + 1 = m(A) \cdot m(c(A)).
\end{equation}

\section{The trivariate generating function}\label{sec:genfn}

\subsection{The functional equation}

\begin{definition}\label{def:genfn}
Define the formal series
\[
G(x,y,z) := \sum_{A \in \SLNz} x^{\chi(A)}\, y^{m(A)}\, z^{m(c(A))} \;\in\; \Z[[x,y,z]].
\]
For each $A \in \SLNz$, $\chi(A), m(A), m(c(A)) \in \Nz$, and the support is locally finite in each variable. The constant term ($A = I$) is $1 \cdot y \cdot z$; we adopt the normalisation $x^0 y^1 z^1 = yz$ for the identity contribution.
\end{definition}

\begin{theorem}\label{thm:fnceq}
The series $G$ satisfies the functional equation
\begin{equation}\label{eq:fnceq}
G(x,y,z) \;=\; yz \;+\; G\bigl(xz^2,\, xyz,\, z\bigr) \;+\; G\bigl(xy^2,\, y,\, xyz\bigr).
\end{equation}
\end{theorem}

\begin{proof}
Every $A \in \SLNz$ either equals $I$ or has a unique left factorisation $A = SA'$ or $A = TA'$ with $A' \in \SLNz$ (free monoid property, \cite[Cor.~1.2]{shakov}). The identity contributes $x^0 y^{1} z^{1} = yz$.

\emph{$S$-children.} For $A = SA'$ with $A' = \bigl(\begin{smallmatrix}a&b\\c&d\end{smallmatrix}\bigr)$, we have $A = \bigl(\begin{smallmatrix}a&b\\a+c&b+d\end{smallmatrix}\bigr)$, hence
\[
m(A) = a^2+b^2 = m(A'), \qquad \chi(A) = a(a+c)+b(b+d) = m(A')+\chi(A'),
\]
and
\[
m(c(A)) = (a+c)^2 + (b+d)^2 = (a^2+b^2) + 2(ac+bd) + (c^2+d^2) = m(A') + 2\chi(A') + m(c(A')).
\]
Thus the $S$-child contribution is
\begin{align*}
\sum_{A'} x^{m(A')+\chi(A')} y^{m(A')} z^{m(A')+2\chi(A')+m(c(A'))}
&= \sum_{A'} (xz^2)^{\chi(A')}\,(xyz)^{m(A')}\, z^{m(c(A'))} \\
&= G(xz^2,\, xyz,\, z).
\end{align*}

\emph{$T$-children.} For $A = TA'$, $A = \bigl(\begin{smallmatrix}a+c&b+d\\c&d\end{smallmatrix}\bigr)$, so
\[
m(c(A)) = c^2+d^2 = m(c(A')), \quad \chi(A) = (a+c)c+(b+d)d = \chi(A') + m(c(A')),
\]
\[
m(A) = (a+c)^2+(b+d)^2 = m(A') + 2\chi(A') + m(c(A')).
\]
Thus the $T$-child contribution is
\begin{align*}
\sum_{A'} x^{\chi(A')+m(c(A'))} y^{m(A')+2\chi(A')+m(c(A'))} z^{m(c(A'))}
&= \sum_{A'} (xy^2)^{\chi(A')}\, y^{m(A')}\, (xyz)^{m(c(A'))} \\
&= G(xy^2,\, y,\, xyz).
\end{align*}
Adding the two and the identity term yields \eqref{eq:fnceq}.
\end{proof}

\begin{remark}
The functional equation \eqref{eq:fnceq} is the exact analogue, for the $\Phzero$-tree, of the Stern--Brocot/Calkin--Wilf functional equation for the rational generating function. It is symmetric under the simultaneous swap $(y,z) \leftrightarrow (z,y)$, which is the manifestation of the complement involution $c$ acting on $\SLNz$ (and which sends $S$-children to $T$-children).
\end{remark}

\subsection{Specialisations}

\begin{corollary}\label{cor:divisor-gf}
Let
\[
\widetilde G(x) := G(x,1,1) = \sum_{A \in \SLNz} x^{\chi(A)} = \sum_{n \ge 0} \tau(n^2+1) \, x^n,
\]
where $\tau(N)$ is the divisor count. Then $\widetilde G$ satisfies
\begin{equation}\label{eq:widetildeG}
\widetilde G(x) = 1 + 2\, G(x, x, 1)\quad \text{(formal identity in $x$),}
\end{equation}
where $G(x,x,1) = \sum_{A} x^{\chi(A)+m(A)}$ counts matrices weighted by the sum of cross-term and top-row norm.
\end{corollary}

\begin{proof}
Set $y = z = 1$ in \eqref{eq:fnceq}: $G(x,1,1) = 1 + G(x, x, 1) + G(x, 1, x)$. By the symmetry $G(x,y,z) = G(x,z,y)$ (which is immediate from the involution $c$ applied to the indexing set: $c$ is a bijection $\SLNz \to \SLNz$ fixing $\chi$ and swapping $m(A) \leftrightarrow m(c(A))$, since $\chi(c(A)) = \chi(A)$ by direct computation $\chi(c(A)) = db + ca = ac+bd = \chi(A)$), the two terms $G(x,x,1)$ and $G(x,1,x)$ are equal. The identity $\sum_A x^{\chi(A)} = \sum_n \tau(n^2+1) x^n$ is Lemma~\ref{lem:tau-count} (the divisor-count theorem of \cite{shakov}).
\end{proof}

\begin{remark}\label{rem:not-rational}
The series $\widetilde G(x) = \sum_n \tau(n^2+1) x^n$ is \emph{not} a rational function of $x$ (its coefficients $\tau(n^2+1)$ have unbounded growth: $\tau(n^2+1) = \Omega((\log n)^{\log 2 - o(1)})$ on average). Hence \eqref{eq:fnceq} should be regarded as an algebraic functional equation in three formal variables, not a polynomial / rational identity.
\end{remark}

\subsection{Gaussian-integer interpretation of the divisor count}

The bijectivity of $\Phzero$ corresponds to the unique factorisation of Gaussian integers $n+i \in \Zi$. We make the correspondence precise.

\begin{proposition}\label{prop:gauss-factor}
For each $A = \bigl(\begin{smallmatrix}a&b\\c&d\end{smallmatrix}\bigr) \in \SLNz$, set $\alpha := a+bi$, $\beta := c+di \in \Zi$. Then
\begin{equation}\label{eq:bar-alpha-beta}
\bar\alpha \cdot \beta = (a-bi)(c+di) = (ac+bd) + i(ad-bc) = \chi(A) + i.
\end{equation}
In particular, $\bar\alpha \beta = n + i$ for $n = \chi(A)$, and the norms satisfy $N(\alpha) = m(A)$, $N(\beta) = m(c(A))$ with $N(\alpha) N(\beta) = N(n+i) = n^2+1$.
\end{proposition}

\begin{proof}
Direct computation: $(a-bi)(c+di) = ac + adi - bci - bdi^2 = (ac+bd) + i(ad-bc)$. Since $A \in \SLNz$ has $ad-bc = 1$, this equals $\chi(A) + i$. Norms are $N(\alpha) = a^2+b^2 = m(A)$ and $N(\beta) = c^2+d^2 = m(c(A))$.
\end{proof}

\begin{corollary}[Gaussian-integer encoding of $\Dphi$]\label{cor:gauss-encoding}
The bijection $\Phzero: \SLNz \to \Dphi$ factors through the Gaussian factorisation map
\[
\SLNz \;\xrightarrow{A \mapsto (\bar\alpha,\beta)}\; \{(\gamma, \delta) \in \Zi \times \Zi : \gamma \delta = n+i \text{ for some } n \in \Nz, \gamma \in \Nz + i \cdot (-\Nz)\} \;\xrightarrow{\text{quotient}}\; \Dphi,
\]
where the second arrow forgets the unit choice $\{\pm 1, \pm i\}$ and records only $(N(\gamma), n)$. Equivalently, $\tau(n^2+1)$ equals the number of unordered factorisations of $n+i$ in $\Zi$ modulo units, which is the count of ideal divisors of $(n+i)$ in $\Zi$.
\end{corollary}

\begin{proof}
The composition $A \mapsto (\bar\alpha, \beta) \mapsto (N(\bar\alpha), \chi(A)) = (m(A), \chi(A)) = \Phzero(A)$ is the original $\Phzero$. Bijectivity \cite[Theorem 2.7]{shakov} is the statement that this composition is a bijection onto $\Dphi$.

The connection to $\tau(n^2+1)$: factorisations of $n+i$ in $\Zi$ (up to units) are in bijection with ideal divisors of the principal ideal $(n+i) \subset \Zi$. Since $\Zi$ is a PID, ideal divisors equal element divisors up to units, and the count is $\tau_{\Zi}((n+i)) = \tau_\Z(N(n+i)) = \tau(n^2+1)$ provided $(n+i)$ is squarefree as an ideal. In general (when $n^2+1$ is not squarefree, e.g., $n=7$ giving $50 = 2\cdot 5^2$), the divisor counts in $\Zi$ and $\Z$ still agree because the rational primes that ramify or split in $\Zi$ are precisely those dividing $n^2+1$ (with the appropriate splitting type), and the divisor function is multiplicative in both $\Z$ and $\Zi$ along these primes.
\end{proof}

\begin{remark}\label{rem:hecke-honest}
A clean Hecke L-function decomposition of the Dirichlet series $\sum_n \tau(n^2+1) n^{-s}$ would proceed by writing $\tau(n^2+1) = \sum_{d \mid n^2+1} 1$ and then decomposing the divisor sum by the splitting type of $d$ in $\Zi$. We do not give a fully clean identity here: the technical complications (non-squarefree $n^2+1$, finite-Euler corrections at the ramified prime $p = 2$) require care that is orthogonal to the spine-sieve discussion. The reader is referred to Cox \cite{cox} for the analytic theory of the relevant Hecke L-functions.
\end{remark}

\section{Sundaram dictionary}\label{sec:sundaram}

We give a precise correspondence between Sundaram's classical sieve and the interior-spine sieve.

\subsection{Sundaram is a degenerate spine sieve}

\begin{theorem}[Sundaram dictionary]\label{thm:sundaram-dict}
Define the linear polynomial $f_S(n) = 2n+1$. The classical Sundaram sieve set
\[
\Sigma := \{i + j + 2ij : i,j \ge 1\} \subset \N
\]
admits the following spine-sieve description. For each $i \ge 1$, set
\[
m_i := 2i+1, \qquad n_i := 3i+1.
\]
Then $(m_i, n_i)$ is a divisor pair for $f_S$ (i.e., $m_i \mid f_S(n_i) = 2(3i+1)+1 = 3(2i+1)$), and
\[
\Sigma \;=\; \bigsqcup_{i \ge 1} \{n_i + k m_i : k \ge 0\} \quad \text{(after deduplication)}.
\]
In particular, Sundaram's sieve is the union of arithmetic progressions $\{(3i+1) + k(2i+1) : k \ge 0\}$ for $i \ge 1$, each anchored at the smallest interior pair with first coordinate $2i+1$.
\end{theorem}

\begin{proof}
For fixed $i \ge 1$, the substitution $j = k+1$ gives $i+j+2ij = i+(k+1)+2i(k+1) = (3i+1)+k(2i+1)$. Hence the parametrisation $(i,j) \mapsto i+j+2ij$ with $i \ge 1, j \ge 1$ enumerates exactly $\bigcup_i \{n_i + km_i : k \ge 0\}$. Each element appears at least once; deduplication (i.e., taking the union of sets rather than the multiset of values) yields $\Sigma$. The verification that $m_i \mid f_S(n_i)$ is direct: $f_S(n_i) = 2(3i+1)+1 = 6i+3 = 3 m_i$.
\end{proof}

\begin{corollary}\label{cor:sundaram-as-spine}
The Sundaram exclusions are the depth-$1$ spine exclusions in the (non-enumerable) divisor-pair set $\mathcal{D}_{f_S}$, where ``anchor'' is parameterised by $\N$ (a single index $i$) rather than by $\SLNz$.
\end{corollary}

\begin{remark}\label{rem:non-enumerable}
The polynomial $f_S(n) = 2n+1$ is not enumerable in the sense of \cite[Defn.~2.1]{shakov}: $|f_S(0)| = 1$ is satisfied, but $f_S$ has degree $1 < 2$, violating one of the size criteria. There is no $\SLNz$-action on $\mathcal{D}_{f_S}$ that makes the divisor pair set bijectively isomorphic to a free monoid; the ``anchor'' index $i$ is just $\N$, the most degenerate possible enumerating set.
\end{remark}

\subsection{Spine sieve carries strictly more information}

\begin{theorem}\label{thm:spine-vs-sundaram}
Consider the analogue of Sundaram applied to $\phi_0(n) = n^2+1$: the naive composite-witness sieve
\[
\Sigma_{\phi_0}^{\mathrm{naive}} := \{n \ge 1 : \exists\, 1 < d < \phi_0(n) \text{ with } d \mid \phi_0(n)\}.
\]
Then $\Sigma_{\phi_0}^{\mathrm{naive}} = \Iset$ (the spine-sieve set of \cite[Defn.~3.4]{B8}), but the spine sieve carries strictly more information in the following sense:
\begin{enumerate}
\item Each $n \in \Iset$ is covered by exactly $\tau(n^2+1) - 2$ interior-anchored spines (counted with multiplicity by which interior pair $(m_0, n_0)$ has $n$ in its AP $\{n_0 + km_0\}$).
\item Each spine carries an algebraic witness $(m_0, n_0)$ derived from a specific Gaussian-integer factorisation of $n_0^2+1$ of norm $m_0$. The naive sieve only knows that $\phi_0(n)$ factors; it does not know \emph{how}.
\end{enumerate}
\end{theorem}

\begin{proof}
\emph{Equality of sets.} If $n \in \Sigma_{\phi_0}^{\mathrm{naive}}$ then $\phi_0(n) = mk$ with $1 < m < \phi_0(n)$, so $(m, n) \in \Dphi$ is a non-boundary (interior) pair. By bijectivity of $\Phzero$, $(m,n) = \Phzero(A)$ for some $A \in \SLN$; thus $n$ is the second component of $\Phzero(A)$ for some interior $A$, hence $n \in \Iset$. Conversely if $n \in \Iset$ then by definition some interior-anchored $A_0$ has $\Phzero(A_0) = (m_0, n_0)$ with $n = n_0+km_0$ for some $k \ge 0$. By Lemma 3.2 of \cite{B8}, $\Phzero(S^k A_0) = (m_0, n)$, giving an interior pair, so $n \in \Sigma_{\phi_0}^{\mathrm{naive}}$.

\emph{Counting (i).} The number of interior pairs $(m,n) \in \Dphi$ with second component $n$ is exactly $\tau(n^2+1) - 2$ (subtract the two boundary pairs $(1,n)$ and $(n^2+1, n)$). Each such interior pair $(m,n)$ corresponds, under $\Phzero^{-1}$, to a unique interior matrix $A \in \SLN$. The matrix $A$ either has its rightmost letter equal to $S$ or to $T$ (or $A$ has only one letter, but then $A$ is on the boundary, contradiction). If the rightmost letter is $S$, write $A = S^k A_0$ with $A_0$ ending in $T$ (or $A_0 = I$); but $A_0 \ne I$ since $A \in \SLN$ means $A$ contains a $T$ as well, so $A_0 \in \SLN$ and is interior-anchored. The pair $\Phzero(A_0) = (m_0, n_0)$ has $m_0 = m$ and $n_0 = n - km$, so $A$ contributes to the $S$-spine of $A_0$ at offset $k$. The $T$-spine analogue is symmetric.

\emph{Algebraic witness (ii).} The matrix $A_0 = \bigl(\begin{smallmatrix}a&b\\c&d\end{smallmatrix}\bigr) \in \SLN$ produces $m_0 = a^2+b^2$ and $n_0 = ac+bd$, with $n_0^2+1 = (a^2+b^2)(c^2+d^2)$ a Gaussian-integer factorisation. The witness records two Gaussian integers $\bar\gamma = a-bi$ and $\delta = c+di$ with $\bar\gamma\delta = (ac+bd)+i(ad-bc) = n_0+i$. Sundaram's analogue would give only $m_0 \mid \phi_0(n_0)$, with no further algebraic structure.
\end{proof}

\begin{remark}\label{rem:gain}
The information gain in (ii) is genuine but \emph{not} obviously analytically usable: the algebraic structure produces additional combinatorial constraints (the witness $(a,b,c,d)$ satisfies $\det = 1$), but these constraints are exactly the Gaussian-prime factorisation that the divisor function already encodes. The hope, articulated in \cite[Sec.~3.4]{B8}, is that the bilinear structure $(a,b) \times (c,d)$ of the witness can feed into a Friedlander--Iwaniec-style Type II estimate; we make no progress on that question here.
\end{remark}

\section{Multiplicity and the divisor function}\label{sec:multiplicity}

\begin{definition}\label{def:mult}
For $n \ge 1$, the \emph{$S$-spine multiplicity} at $n$ is
\[
\mu^{(S)}_{\phi_0}(n) := \#\{ A_0 \in \SLN : A_0 \text{ ends in } T,\ \chi(A_0) \le n,\ m(A_0) \mid n - \chi(A_0)\},
\]
and the \emph{$T$-spine multiplicity} $\mu^{(T)}_{\phi_0}(n)$ is defined symmetrically by swapping $S \leftrightarrow T$ throughout.
\end{definition}

\begin{theorem}\label{thm:multiplicity}
For all $n \ge 1$,
\[
\mu^{(S)}_{\phi_0}(n) + \mu^{(T)}_{\phi_0}(n) = \tau(n^2+1) - 2.
\]
\end{theorem}

\begin{proof}
Each interior matrix $A \in \SLN$ has a unique rightmost letter, $S$ or $T$. If $A$ ends in $T$, then $A = S^0 \cdot A$ realises $A$ as an $S$-spine matrix above its anchor $A_0 = A$ (which itself ends in $T$); more generally, any $A$ ending in $T$ is reached as $S^k A_0$ for $A_0 := A$ ($k = 0$). Conversely, $A = S^k A_0$ for some $k \ge 0$ and $A_0$ ending in $T$ is a unique decomposition (by the free monoid property of $\SLNz$, peeling off the maximal $S$-prefix). Thus
\[
\#\{A \in \SLN : \chi(A) = n,\ A \text{ ends in } T\} = \#\{(A_0, k) : A_0 \in \SLN \text{ ends in } T, k \ge 0, \chi(S^k A_0) = n\}.
\]
The condition $\chi(S^k A_0) = n$ is $\chi(A_0) + k \cdot m(A_0) = n$, which has a unique nonneg solution $k = (n - \chi(A_0))/m(A_0)$ if and only if $\chi(A_0) \le n$ and $m(A_0) \mid (n - \chi(A_0))$. Thus the right-hand side equals $\mu^{(S)}_{\phi_0}(n)$.

The total count $\mu^{(S)}_{\phi_0}(n) + \mu^{(T)}_{\phi_0}(n)$ therefore equals $\#\{A \in \SLN : \chi(A) = n\}$. By Lemma \ref{lem:tau-count} below, this is $\tau(n^2+1) - 2$.
\end{proof}

\begin{lemma}[from \cite{shakov}]\label{lem:tau-count}
For $n \ge 0$, $\#\{A \in \SLNz : \chi(A) = n\} = \tau(n^2+1)$. The two boundary matrices $A = S^n$ and $A = T^n$ are the unique non-interior matrices with $\chi = n$, hence
\[
\#\{A \in \SLN : \chi(A) = n\} = \tau(n^2+1) - 2 \quad \text{for}\ n \ge 1.
\]
\end{lemma}

\begin{proof}
The first identity is the divisor-counting theorem of \cite{shakov} (labelled \texttt{divisorcount} therein). For the boundary count: $\chi(S^\alpha) = 1 \cdot \alpha + 0 \cdot 1 = \alpha$ from $S^\alpha = \bigl(\begin{smallmatrix}1 & 0 \\ \alpha & 1\end{smallmatrix}\bigr)$, and similarly $\chi(T^\alpha) = 0 + \alpha = \alpha$. So for $\chi = n$ the boundary contributes exactly $\{S^n, T^n\}$, two matrices.
\end{proof}

\begin{corollary}\label{cor:mult-prime}
$\phi_0(n)$ is prime $\iff \mu^{(S)}_{\phi_0}(n) + \mu^{(T)}_{\phi_0}(n) = 0$. \end{corollary}

This recasts Landau's fourth problem as the following purely combinatorial statement.

\begin{conjecture}[Landau IV, restated]\label{conj:landau-mult}
For infinitely many $n \in \N$, no $A_0 \in \SLN$ produces a pair $\Phzero(A_0) = (m_0, n_0)$ with $n \equiv n_0 \pmod{m_0}$ and $n_0 \le n$.
\end{conjecture}

\section{Depth-bounded densities}\label{sec:depth}

For $A \in \SLN$, define $\depth(A)$ to be the $S/T$-word length of $A$. The boundary matrices ($I$, $S^\alpha$, $T^\alpha$ for $\alpha \ge 1$) have depth $0$ or $1$ but are not interior. The smallest interior depth is $2$ ($ST$ and $TS$).

\begin{definition}\label{def:depth-spines}
For $k \ge 2$, let
\[
\Iset^{(\le k)} := \bigcup_{\substack{A_0 \in \SLN \\ \depth(A_0) \le k}} \{n_0(A_0) + j m_0(A_0) : j \ge 0\}, \qquad (m_0, n_0) := \Phzero(A_0),
\]
and let $\mathcal{P}^{(\le k)}_N := \{n \in [1,N] : n \notin \Iset^{(\le k)}\}$ be the set of \emph{depth-$\le k$ survivors} up to $N$.
\end{definition}

\begin{remark}\label{rem:overcount}
The union $\Iset^{(\le k)}$ overcounts: a fixed $n$ may lie in many depth-$\le k$ spines simultaneously. The complement $\mathcal{P}^{(\le k)}_N$ is straightforward to compute via inclusion--exclusion on a finite set of APs.
\end{remark}

\subsection{Exact density at depth $\le 2$}

The interior matrices of depth exactly $2$ are precisely $ST$ and $TS$. Compute:
\[
\Phzero(ST) = \Phzero\!\!\begin{pmatrix}1&1\\1&2\end{pmatrix} = (2, 3), \qquad \Phzero(TS) = \Phzero\!\!\begin{pmatrix}2&1\\1&1\end{pmatrix} = (5, 3).
\]
The associated APs are $\{3+2j : j \ge 0\} = \{3,5,7,9,\dots\}$ (odd $n \ge 3$) and $\{3+5j : j \ge 0\} = \{3,8,13,18,\dots\}$ (i.e., $n \equiv 3 \pmod 5$, $n \ge 3$).

\begin{theorem}\label{thm:depth2-density}
The set of depth-$\le 2$ survivors $\mathcal{P}^{(\le 2)}_N$ has density
\[
\lim_{N \to \infty} \frac{|\mathcal{P}^{(\le 2)}_N|}{N} = \frac{2}{5}.
\]
Concretely: $n$ is a depth-$\le 2$ survivor iff $n \in \{1, 2\}$ or ($n \ge 4$, $n$ even, and $n \not\equiv 3 \pmod 5$).
\end{theorem}

\begin{proof}
The two depth-$2$ exclusions cover odd $n \ge 3$ and $n \equiv 3 \pmod 5$ ($n \ge 3$). Their complement among $n \ge 4$ consists of even $n$ with $n \not\equiv 3 \pmod 5$. Even $n$ have density $1/2$; among even $n$, the residue class $n \equiv 8 \pmod{10}$ (the unique residue mod $10$ that is even and $\equiv 3 \pmod 5$) has density $1/10$ in $\N$. Surviving density: $1/2 - 1/10 = 4/10 = 2/5$.

The two extra survivors $n=1,2$ satisfy $\phi_0(1)=2$, $\phi_0(2)=5$, both prime, consistent with depth-$2$ giving no exclusion in the range $[1,2]$.
\end{proof}

\begin{remark}
The depth-$2$ density $2/5$ is far above any heuristic prediction for primes among $\{n^2+1\}$: Hardy--Littlewood predicts $\pi_{n^2+1}(N) \sim 1.3727 \cdot N/\log N$, hence density tending to $0$. The depth-$2$ survivors are dominated by composites of $\phi_0$ that have a divisor not detected by the two depth-$2$ spines.
\end{remark}

\subsection{Depth $\le 3$ exact density}

The interior matrices of depth $3$ (with both letters present): $S^2T, ST^2, T^2S, TS^2, STS, TST$. Compute:

\[
\begin{array}{lcll}
A & = & \Phzero(A) & \text{spine AP modulo $m_0$} \\
\hline
S^2T & \begin{pmatrix}1&1\\2&3\end{pmatrix} & (2, 5) & n \equiv 1 \pmod 2 \\
ST^2 & \begin{pmatrix}1&2\\1&3\end{pmatrix} & (5, 7) & n \equiv 2 \pmod 5 \\
T^2S & \begin{pmatrix}3&2\\1&1\end{pmatrix} & (13, 5) & n \equiv 5 \pmod{13} \\
TS^2 & \begin{pmatrix}3&1\\2&1\end{pmatrix} & (10, 7) & n \equiv 7 \pmod{10} \\
STS & \begin{pmatrix}2&1\\3&2\end{pmatrix} & (5, 8) & n \equiv 3 \pmod 5 \\
TST & \begin{pmatrix}2&3\\1&2\end{pmatrix} & (13, 8) & n \equiv 8 \pmod{13}
\end{array}
\]

\begin{lemma}\label{lem:depth3-redundancy}
Among the six depth-$3$ spines, the spines from $S^2T$ and $TS^2$ are sub-progressions of depth-$\le 2$ spines:
\begin{itemize}
\item $S^2T$: $n \equiv 1 \pmod 2$ is exactly the depth-$2$ spine from $ST$ (which gave $n \equiv 1 \pmod 2$ from $n \ge 3$).
\item $TS^2$: $n \equiv 7 \pmod{10}$, in particular $n$ odd, is contained in the $ST$-spine.
\item $STS$: $n \equiv 3 \pmod 5$, in particular contained in the $TS$-spine $\{n \equiv 3 \pmod 5\}$.
\end{itemize}
The non-redundant new constraints from depth $3$ are: $n \equiv 2 \pmod 5$ (from $ST^2$), $n \equiv 5 \pmod{13}$ (from $T^2S$), and $n \equiv 8 \pmod{13}$ (from $TST$).
\end{lemma}

\begin{proof}
Direct verification of containments by computing residues: $1 \pmod 2 \subseteq 1 \pmod 2$; $\{7,17,27,\dots\} \pmod{10}$ all odd, hence in the depth-$2$ odd-spine; $\{3,8,13,18,\dots\} \pmod 5 = \{3 \pmod 5\}$ matches.
\end{proof}

\begin{theorem}\label{thm:depth3-density}
\[
\lim_{N \to \infty} \frac{|\mathcal{P}^{(\le 3)}_N|}{N} = \frac{33}{130}.
\]
\end{theorem}

\begin{proof}
By Lemma \ref{lem:depth3-redundancy}, the union of depth-$\le 3$ exclusions reduces to the five arithmetic progressions
\[
n \equiv 1 \pmod 2,\quad n \equiv 3 \pmod 5,\quad n \equiv 2 \pmod 5,\quad n \equiv 5 \pmod{13},\quad n \equiv 8 \pmod{13}.
\]
We compute the surviving density by passing to residues modulo $\mathrm{lcm}(2,5,13) = 130$.

Step 1: $n$ even (the complement of the first AP) gives $65$ residues mod $130$.

Step 2: among even $n$, exclude $n \equiv 2 \pmod 5$ and $n \equiv 3 \pmod 5$. Among even residues mod $10$, namely $\{0,2,4,6,8\}$, the residues $2 \pmod 5$ and $3 \pmod 5$ are realised by $\{2\}$ and $\{8\}$ respectively. The surviving even residues mod $10$ are $\{0,4,6\}$. By CRT modulo $130$ this gives $3 \cdot 13 = 39$ surviving residues.

Step 3: from these $39$, exclude $n \equiv 5 \pmod{13}$ and $n \equiv 8 \pmod{13}$. By $\gcd(10,13)=1$ and CRT, exactly $2/13$ of the $39$ surviving residues have $n \equiv 5$ or $8 \pmod {13}$, i.e., $39 \cdot 2/13 = 6$ residues mod $130$.

Surviving count: $39 - 6 = 33$ residues mod $130$, hence density $33/130$.
\end{proof}

\begin{remark}
Direct enumeration of $n \in [1, 130]$ confirms the count. The $33$ surviving values are
\[
\{4, 6, 10, 14, 16, 20, 24, 26, 30, 36, 40, 46, 50, 54, 56, 64, 66, 74, 76, 80, 84, 90, 94, 100, 104, 106, 110, 114, 116, 120, 124, 126, 130\}.
\]
Checking that $\phi_0$ at each survivor has no divisor witnessed by a depth-$\le 3$ matrix is automatic from the construction. Checking primality of $\phi_0$ at each survivor is a separate question; e.g., $\phi_0(4) = 17$ prime, $\phi_0(6) = 37$ prime, $\phi_0(10) = 101$ prime, $\phi_0(14) = 197$ prime, but $\phi_0(16) = 257$ prime, $\phi_0(20) = 401$ prime, $\phi_0(24) = 577$ prime, $\phi_0(26) = 677$ prime, $\phi_0(30) = 901 = 17 \cdot 53$ \emph{composite} (witnessed by a depth-$> 3$ spine of modulus $17$). So depth-$\le 3$ survival is genuinely weaker than primality, as required.
\end{remark}

\begin{corollary}\label{cor:depth3-survivors}
$|\mathcal{P}^{(\le 3)}_N| = (33/130) N + O(1)$. In particular, an asymptotic positive density of $n$ survives all spine exclusions of depth $\le 3$.
\end{corollary}

\subsection{Asymptotic surviving density at depth $\le k$}

\begin{theorem}\label{thm:depth-k-upper}
For each $k \ge 2$, the depth-$\le k$ surviving density
\[
\delta_k := \lim_{N \to \infty} \frac{|\mathcal{P}^{(\le k)}_N|}{N}
\]
exists, satisfies $\delta_2 \ge \delta_3 \ge \delta_4 \ge \cdots$, and is bounded below by the singular series ratio
\begin{equation}\label{eq:singular-series-bound}
\delta_k \;\ge\; \prod_{p \le M(k)} \!\left(1 - \frac{\rho_{\phi_0}(p)}{p}\right),
\end{equation}
where $\rho_{\phi_0}(p) := \#\{n \pmod p : p \mid n^2+1\}$ and $M(k)$ is the largest prime $p$ such that some interior matrix $A_0$ of depth $\le k$ has $m(A_0) = p$.
\end{theorem}

\begin{proof}
The depth-$\le k$ exclusions form a finite union of arithmetic progressions. The density of the complement of a finite union of APs in $\N$ equals the natural density of its CRT-residue classes, which exists. The monotonicity $\delta_{k+1} \le \delta_k$ is from $\Iset^{(\le k)} \subseteq \Iset^{(\le k+1)}$.

For the lower bound: the spine APs at depth $\le k$ have moduli $m_0 = a^2+b^2$ for matrices $A_0 = \bigl(\begin{smallmatrix}a&b\\c&d\end{smallmatrix}\bigr)$ of depth $\le k$. The set of $m_0$-values that appear is contained in $\{m \in \N : m = a^2+b^2, a,b \ge 1\}$. For each such $m_0$, the spine AP $\{n_0+jm_0\}$ covers $1$ residue class mod $m_0$. Each \emph{prime} $p$ appearing as $m_0$ contributes $\rho_{\phi_0}(p)$ excluded residue classes mod $p$ (since $p$ divides $n^2+1$ iff $n$ is one of $\rho_{\phi_0}(p)$ residues, and each is hit by a depth-$\le k$ spine of modulus $p$ provided one such spine exists).

For composite $m_0$, the AP excludes at most one residue class. Hence the density of excluded $n$ at depth $\le k$ is at most the density of $n$ excluded by the prime-modulus spines alone, i.e.,
\[
1 - \delta_k \;\le\; 1 - \prod_{\substack{p \le M(k) \\ p\ \text{a depth-$\le k$ spine modulus}}} \!\!\left(1 - \frac{\rho_{\phi_0}(p)}{p}\right).
\]
Since at depth $\le k$ all primes $p \le k$ that are sums of two squares with both summands $\ge 1$ appear as moduli (e.g., $5 = 1+4 = m(ST^2)$, etc.), the product is over all primes $p \le M(k)$ with $\rho_{\phi_0}(p) \ge 1$.
\end{proof}

\begin{remark}\label{rem:singular-series-evaluation}
The product $\prod_{p \le M(k)} (1 - \rho_{\phi_0}(p)/p)$ is the local part of the Hardy--Littlewood singular series for $\{n^2+1\}$. For $\phi_0$, $\rho_{\phi_0}(p) = 2$ if $p \equiv 1 \pmod 4$, $=1$ if $p = 2$, $=0$ if $p \equiv 3 \pmod 4$. As $M(k) \to \infty$, the product tends to $0$ (it diverges like $1/(\log M)$ by Mertens-type arguments restricted to primes $\equiv 1 \pmod 4$). Hence Theorem \ref{thm:depth-k-upper} alone gives $\delta_k \to 0$, but only at rate $1/\log M(k)$, which is the trivial sieve rate.
\end{remark}

\begin{conjecture}[Asymptotic depth survival]\label{conj:depth-asymp}
There exist constants $c_1, c_2 > 0$ and a function $M(k) \to \infty$ such that
\[
\delta_k = \frac{c_1 + o(1)}{\log M(k)} \quad \text{as}\ k \to \infty,
\]
and $M(k) \asymp \lambda_+^k$ where $\lambda_+ = (5+\sqrt{17})/2$ is the dominant eigenvalue of the row-sum recursion of \cite[Cor.~ closed\_forms]{shakov}. Equivalently, $\delta_k \asymp 1/k$ as $k \to \infty$.
\end{conjecture}

\begin{remark}
Conjecture \ref{conj:depth-asymp} is consistent with Hardy--Littlewood for $\{n^2+1\}$: the depth-$\to \infty$ limit gives $\delta_\infty = 0$, but at depth $k$ the survivors include all $n \le N$ such that $n^2+1$ has no prime factor up to $M(k)$, which by sieve heuristics has density $\sim 1/\log M(k) \sim 1/k$. We do not prove the conjecture; the upper bound from Theorem \ref{thm:depth-k-upper} matches the conjecture, but the lower bound (which would assert that ``most'' interior matrices of depth $\le k$ have small modulus) requires a Heuberger--Krenn-type asymptotic result for the level-set distribution of $m(A_0)$ across depth-$k$ matrices.
\end{remark}

\section{Second-tier extension of the generalised lemma}\label{sec:second-tier}

The generalised lemma of \cite[Theorem 1.3]{B8} states that $|f(n)|$ is prime for $n \in [1, |f(1)|]$. We give a sharp second-tier criterion that extends primality to a specific residue class modulo $|f(1)|$.

\begin{theorem}[Second-tier criterion]\label{thm:second-tier}
Let $f \in \{\phi_0, \phi_1, \psi_2, \phi_3\}$ be enumerable, and write $\beta := |f(1)|$. The divisibility relation $\beta \mid f(n)$ holds if and only if $n \equiv 1 \pmod \beta$ (in the cases $f = \phi_0, \phi_1, \phi_3$) or $n \equiv 1 \pmod 2$ (in the case $f = \psi_2$). Furthermore, defining $\beta'$ as the smallest prime $> \beta$ such that $\beta' \mid f(n)$ for some $n \in \N$, we have
\[
\beta' = 5,\ 7,\ 7,\ 11 \quad \text{for } f = \phi_0,\phi_1,\psi_2,\phi_3 \text{ respectively}.
\]
Hence: for each $f$, no prime $p \in (\beta, \beta')$ divides $f(n)$ for any $n$.
\end{theorem}

\begin{proof}
\emph{Divisibility by $\beta$.}
For $\phi_0$: $\phi_0(n) = n^2+1$. We have $\phi_0(n) \equiv 0 \pmod 2$ iff $n$ is odd. Here $\beta = 2$, and the criterion is exactly the parity statement.

For $\phi_1$: $\phi_1(n) = n^2+n+1$, $\beta = 3$. Compute mod 3: $n \equiv 0: 1$; $n \equiv 1: 0$; $n \equiv 2: 1$. So $3 \mid \phi_1(n)$ iff $n \equiv 1 \pmod 3$.

For $\psi_2$: $\psi_2(n) = n^2+2n-1$, $\beta = 2$. Mod 2: $n^2 - 1$, which vanishes iff $n$ odd.

For $\phi_3$: $\phi_3(n) = n^2+3n+1$, $\beta = 5$. Mod 5: $n \equiv 0: 1$; $n \equiv 1: 0$; $n \equiv 2: 1$; $n \equiv 3: 4$; $n \equiv 4: 4$. So $5 \mid \phi_3(n)$ iff $n \equiv 1 \pmod 5$.

\emph{Computation of $\beta'$.}
For each of the four polynomials, we verify by direct mod-$p$ computation that no prime $p \in (\beta, \beta')$ divides $f(n)$ for any $n$.

$\phi_0$, $\beta=2,\beta'=5$: $\phi_0(n) \pmod 3 = n^2+1 \in \{1, 2\}$, never $0$.

$\phi_1$, $\beta=3, \beta'=7$: $\phi_1(n) \pmod 5 \in \{1, 3, 2, 3, 1\}$ for $n \equiv 0,\ldots,4$, never $0$.

$\psi_2$, $\beta=2,\beta'=7$: $\psi_2(n) \pmod 3 = n^2+2n-1 \in \{2, 2, 1\}$ for $n \equiv 0,1,2$, never $0$. $\psi_2(n) \pmod 5 \in \{4, 2, 2, 4, 3\}$ for $n \equiv 0,\ldots,4$, never $0$.

$\phi_3$, $\beta=5, \beta'=11$: $\phi_3(n) \pmod 7 = n^2+3n+1 \in \{1, 5, 4, 5, 1, 6, 6\}$ for $n \equiv 0,\ldots,6$, never $0$.

That $\beta'$ does divide $f$ at some $n$ is verified by the explicit witnesses $\phi_0(7) = 50 = 2 \cdot 5^2$ (so $5 \mid \phi_0(7)$); $\phi_1(2) = 7$; $\psi_2(2) = 7$; $\phi_3(2) = 11$.
\end{proof}

\begin{corollary}[Second-tier interval]\label{cor:second-tier-interval}
For each enumerable $f$, in the second-tier interval $\{n \in [\beta+1,\beta'-1] : n \not\equiv 1 \pmod \beta \text{ (or $n$ even for $\psi_2$)}\}$, every interior pair $(m, n) \in \Df$ has $m \ge \beta'$. In particular, no depth-$\le k$ spine of modulus $\le \beta$ excludes any such $n$ for any $k$.
\end{corollary}

\begin{proof}
By Theorem \ref{thm:second-tier}, $\beta \nmid f(n)$ for $n$ in this interval, and no prime $p \in (\beta, \beta')$ divides $f(n)$. Hence any nontrivial divisor $m$ of $f(n)$ has all prime factors $\ge \beta'$, so $m \ge \beta'$.
\end{proof}

\begin{remark}\label{rem:second-tier-empirical}
Corollary \ref{cor:second-tier-interval} is not the statement that $f(n)$ is prime in the second-tier interval: $f(n)$ might still be divisible by some $m \ge \beta'$. Empirically:
\begin{itemize}
\item $\phi_0$: $\beta = 2, \beta' = 5$. Interval $\{n \in [3, 4], n$ even$\} = \{4\}$. $\phi_0(4) = 17$ prime. eckmark
\item $\phi_1$: $\beta = 3, \beta' = 7$. Interval $\{n \in [4,6], n \not\equiv 1 \pmod 3\} = \{5, 6\}$ (excluding $n = 4 \equiv 1$ but $4 \in [4,6]$ should be checked: $4 \equiv 1 \pmod 3$, so excluded). Remaining $\{5, 6\}$. $\phi_1(5) = 31$ prime; $\phi_1(6) = 43$ prime. eckmarkeckmark
\item $\psi_2$: $\beta = 2, \beta' = 7$. Interval $\{n \in [3,6], n$ even$\} = \{4, 6\}$. $\psi_2(4) = 23$ prime; $\psi_2(6) = 47$ prime. eckmarkeckmark
\item $\phi_3$: $\beta = 5, \beta' = 11$. Interval $\{n \in [6, 10], n \not\equiv 1 \pmod 5\} = \{7, 8, 9, 10\}$. $\phi_3(7) = 71$ prime; $\phi_3(8) = 89$ prime; $\phi_3(9) = 109$ prime; $\phi_3(10) = 131$ prime. eckmarkeckmarkeckmarkeckmark
\end{itemize}
The empirical primality is striking but \emph{not} proven by Theorem \ref{thm:second-tier}: the theorem only excludes the smallest possible interior witness. Larger interior witnesses with $m \in [\beta', f(n)/\beta')$ remain possible; that none of them appear empirically reflects the small height of $f(n)$ in the second-tier interval.
\end{remark}

\begin{question}\label{q:second-tier}
Is it true that for every enumerable $f$ and every $n$ in the second-tier interval $[\beta+1, \beta'-1]$ with $n \not\equiv 1 \pmod \beta$ (resp.\ $n$ even for $\psi_2$), $|f(n)|$ is prime?
\end{question}

\begin{remark}
A positive answer to Question \ref{q:second-tier} would extend the a-priori prime count of \cite[Example 1.4]{B8} substantially. We have empirically verified the question for the four polynomials in their respective second-tier intervals (above). A proof would presumably need a finer analysis of all interior matrices $A_0$ with $m(A_0) \in [\beta', \sqrt{f(\beta'-1)}]$, paralleling the proof of the original generalised lemma.
\end{remark}

\section{Coupled-spine sieve and the highway}\label{sec:joint}

The four enumerable polynomials are coupled through the underlying $\SLNz$: each $A \in \SLNz$ produces simultaneously a divisor pair in each $\mathcal{D}_f$ for $f \in \{\phi_0, \phi_1, \psi_2, \phi_3\}$.

\begin{lemma}\label{lem:joint-cross}
For $A = \bigl(\begin{smallmatrix}a&b\\c&d\end{smallmatrix}\bigr) \in \SLN$, define the four cross-terms
\[
\chi_0 := ac+bd, \quad \chi_1 := ac+bc+bd, \quad \chi_2 := \max(ac,bd) + 2bc - \min(ac,bd), \quad \chi_3 := ac+3bc+bd.
\]
Then
\[
\chi_1 - \chi_0 = bc, \qquad \chi_3 - \chi_0 = 3 bc, \qquad \chi_2 - \chi_0 = 2bc - 2\min(ac, bd).
\]
On $\SLN$ where $bc \ge 1$, this gives the chain $\chi_0 < \chi_1 < \chi_3$ with $\chi_3 - \chi_1 = 2bc \ge 2$.
\end{lemma}

\begin{proof}
Direct algebra: $\chi_1 = ac+bc+bd = \chi_0 + bc$ and $\chi_3 = ac+3bc+bd = \chi_0+3bc$. For $\chi_2$, split on the sign of $ac - bd$. If $ac \ge bd$, $\chi_2 = ac+2bc-bd$, so $\chi_2 - \chi_0 = 2bc - 2bd = 2bc - 2\min(ac,bd)$. If $ac < bd$, $\chi_2 = bd+2bc-ac$, so $\chi_2 - \chi_0 = 2bc - 2ac = 2bc - 2\min(ac,bd)$. The chain is immediate.
\end{proof}

\begin{corollary}\label{cor:phi0-phi3-coupling}
Every interior $A \in \SLN$ contributes a depth-aligned pair of spine exclusions: an exclusion of $n_0(A) = \chi_0(A) = ac+bd$ from the $\phi_0$-prime set, and an exclusion of $n_3(A) = \chi_3(A) = \chi_0(A) + 3bc$ from the $\phi_3$-prime set. The map $n_0 \mapsto n_3$ is injective on a fixed depth level when $bc$ varies with depth.
\end{corollary}

\begin{conjecture}[Joint density]\label{conj:joint-density}
Define
\[
\Pi^{\mathrm{joint}}(N) := \#\{n \in [1,N] : f(n)\ \text{prime for some}\ f \in \{\phi_0, \phi_1, \psi_2, \phi_3\}\}.
\]
Then $\Pi^{\mathrm{joint}}(N) \to \infty$ as $N \to \infty$.
\end{conjecture}

\begin{remark}
Conjecture \ref{conj:joint-density} is strictly weaker than each individual Bunyakovsky conjecture for the four polynomials, so a-priori it is more accessible. However, the four cross-terms are tightly correlated (Lemma \ref{lem:joint-cross}), and the events ``$f_i(n)$ prime'' are not independent. We make no progress on Conjecture \ref{conj:joint-density} here. A serious attempt would presumably need Lemma \ref{lem:joint-cross} packaged as a coupled spine-sieve estimate; we note that even the straightforward depth-$\le 2$ joint exclusion analysis gives a positive density of $n$ for which \emph{at least one} of $\phi_0(n), \phi_1(n), \psi_2(n), \phi_3(n)$ avoids the smallest depth-$2$ spines, recovering Conjecture \ref{conj:joint-density} only in the trivial sense that this density is positive for finite $N$.
\end{remark}

\section{Status and remaining gaps}\label{sec:status}

We summarise what is rigorously proven, what is conjectured, and what gaps remain to Landau's fourth problem.

\textbf{Proven theorems.}
\begin{itemize}
\item Theorem \ref{thm:fnceq}: Exact functional equation for the trivariate generating function $G(x,y,z)$ of $\Phzero$, derived from the free monoid structure.
\item Corollary \ref{cor:divisor-gf}: Specialisation $\widetilde G(x) = \sum_n \tau(n^2+1) x^n = 1 + 2 G(x, x, 1)$.
\item Proposition \ref{prop:gauss-factor} and Corollary \ref{cor:gauss-encoding}: Explicit Gaussian-integer encoding $A \mapsto (\bar\alpha, \beta)$ realising $\bar\alpha\beta = \chi(A)+i$, and the resulting bijection $\Dphi \leftrightarrow $ unordered Gaussian factorisations of $n+i$.
\item Theorem \ref{thm:sundaram-dict}: Sundaram's classical sieve is exactly the depth-$1$ spine sieve of the (non-enumerable) polynomial $f_S(n) = 2n+1$ in the degenerate enumeration $\N$.
\item Theorem \ref{thm:spine-vs-sundaram}: The spine sieve of $\phi_0$ coincides set-wise with the naive composite-witness sieve, but with strictly more multiplicity / algebraic structure (Theorem \ref{thm:multiplicity}).
\item Theorem \ref{thm:multiplicity}: Multiplicity-counting identity, $\mu^{(S)}_{\phi_0}(n) + \mu^{(T)}_{\phi_0}(n) = \tau(n^2+1) - 2$.
\item Theorem \ref{thm:depth2-density}: The depth-$\le 2$ surviving density is exactly $2/5$.
\item Theorem \ref{thm:depth3-density}: The depth-$\le 3$ surviving density is exactly $33/130$.
\item Theorem \ref{thm:depth-k-upper}: Surviving density at depth $\le k$ is bounded below by the truncated singular series (with truncation tied to $M(k)$).
\item Theorem \ref{thm:second-tier}: Sharp characterisation of when the smallest divisor $\beta = |f(1)|$ divides $f(n)$, giving a second-tier prime-detection interval.
\item Lemma \ref{lem:joint-cross}: Algebraic relations among the four cross-terms $\chi_0, \chi_1, \chi_2, \chi_3$.
\end{itemize}

\textbf{Conjectures.}
\begin{itemize}
\item Conjecture \ref{conj:landau-mult}: A combinatorial restatement of Landau IV.
\item Conjecture \ref{conj:depth-asymp}: $\delta_k \asymp 1/k$ asymptotically; consistent with Hardy--Littlewood.
\item Conjecture \ref{conj:joint-density}: Joint Bunyakovsky for the four polynomials.
\end{itemize}

\textbf{Heuristics / questions.}
\begin{itemize}
\item Question \ref{q:second-tier}: Does $|f(n)|$ remain prime in the entire second-tier interval $[\beta+1, \beta'-1]$ avoiding the residue $1 \pmod \beta$?
\end{itemize}

\textbf{Remaining gaps to Landau IV.} The combinatorial reformulation of Landau IV in terms of spine-sieve survival (Conjecture \ref{conj:landau-mult}) is equivalent to the original. The depth-bounded density results (Theorems \ref{thm:depth2-density}, \ref{thm:depth3-density}, \ref{thm:depth-k-upper}) show that the spine sieve produces the correct \emph{type} of singular-series contribution for $\phi_0$, but the bound $\delta_k \to 0$ at rate $1/\log M(k)$ matches the trivial sieve and does not break parity. The \emph{essential analytic input missing} is a Type II / bilinear estimate of the form posed in \cite[Question 4.1]{B8}:
\[
\sum_{\substack{A_0 \in \SLNz \\ \chi(A_0) \le N}} \alpha(a,b) \beta(c,d) \;=\; o(N) \quad \text{(bounded sequences $\alpha, \beta$)}.
\]
The functional equation \eqref{eq:fnceq} packages the combinatorial input but does not in itself provide bilinear cancellation: the variables $y, z$ track row-norms but not the multiplicative structure modulo small primes, which is what bilinear estimates exploit.

\textbf{Honest verdict.} This paper sharpens the combinatorial side of the spine-sieve framework — exact generating-function equation, exact small-depth densities, exact dictionary with Sundaram, sharp second-tier criterion — but does not advance on the analytic obstacles to Landau IV. The most promising direction we identify is the empirical strength of the second-tier criterion (Question \ref{q:second-tier}), which is a finite-checkable conjecture for each enumerable polynomial and which, if proved, would extend the a-priori prime count from $|f(1)|$ to a substantially larger interval governed by $\beta' = $ next admissible modulus. We have verified Question \ref{q:second-tier} for $\phi_0, \phi_1, \psi_2, \phi_3$ in the relevant intervals.

\begin{thebibliography}{99}

\bibitem{shakov}
A. Shakov, ``Polynomials in $\Z[x]$ whose divisors are enumerated by $\SLNz$,'' February 2024. arXiv:2405.03552.

\bibitem{B8}
(working note) ``A-priori primes for enumerable polynomials and the interior-spine sieve,'' B8 in the Shakov bridge series.

\bibitem{sundaram}
S. P. Sundaram, ``A new sieve for prime numbers,'' \emph{Math. Student} \textbf{2} (1934), 73.

\bibitem{allouche-shallit}
J.-P. Allouche and J. Shallit, ``The ring of $k$-regular sequences,'' \emph{Theor. Comput. Sci.} \textbf{98} (1992), 163--197; \emph{Automatic Sequences}, Cambridge University Press, 2003.

\bibitem{heuberger-krenn}
C. Heuberger and D. Krenn, ``Asymptotic analysis of regular sequences,'' \emph{Algorithmica} \textbf{82} (2020), 429--508. arXiv:1810.13178.

\bibitem{northshield}
S. Northshield, ``Stern's diatomic sequence $0, 1, 1, 2, 1, 3, 2, 3, 1, 4, \ldots$,'' \emph{Amer. Math. Monthly} \textbf{117} (2010), 581--598.

\bibitem{iwaniec1978}
H. Iwaniec, ``Almost-primes represented by quadratic polynomials,'' \emph{Invent. Math.} \textbf{47} (1978), 171--188.

\bibitem{FI-asymptotic}
J. Friedlander and H. Iwaniec, ``Asymptotic sieve for primes,'' \emph{Ann. of Math.} \textbf{148} (1998), 1041--1065.

\bibitem{cox}
D. A. Cox, \emph{Primes of the Form $x^2 + ny^2$}, Wiley, 1989; AMS Chelsea, 3rd ed. 2022.

\bibitem{calkin-wilf}
N. Calkin and H. S. Wilf, ``Recounting the rationals,'' \emph{Amer. Math. Monthly} \textbf{107} (2000), 360--363.

\end{thebibliography}

\end{document}
